\expandafter\def\csname HistoricalNoteText10\endcsname{Satis bene certe demonstratum esse videtur, curvam algebraicam neque alicubi subito abrumpi posse (uti e. g. evenit in curva transscendente, cuius aequatio $y=\frac1{\mathrm{log}\,x}$), neque post spiras infinitas in aliquo puncto se quasi perdere (ut spiralis logarithmica), quantumque scio nemo dubium contra hanc rem movit. Attamen si quis postulat, demonstrationem nullis dubiis obnoxiam alia occasione tradere suscipiam. In casu praesenti vero manifestum est, si aliquis ramus e. g. 2, ex circulo nullibi exiret (fig. 3), te in circulum inter 0 et 2 intrare, postea circa totum hunc ramum (qui in circuli spatio se perdere deberet) circummeare, et tandem inter 2 et 4 rursus ex circulo egredi posse, ita ut nullibi in tota via in lineam primam incideris. Hoc vero absurdum esse inde patet, quod in puncto, ubi in circulum ingressus es, superficiem primam supra te habuisti, in egressu, infra; quare necessario alicubi in superficiem primam ipsam incidere debuisti, sive in punctum lineae primae. --- Ceterum ex hoc ratiocinio principiis geometriae situs innixo, quae haud minus valida sunt, quam principia geometriae magnitudinis, sequitur tantummodo, si in aliquo ramo lineae primae in circulum intres, te alio loco ex circulo rursus egredi posse, semper in linea prima manendo, neque vero, viam tuam esse lineam continuam in eo sensu, quo in geometria sublimiori accipitur. Sed hic sufficit, viam esse lineam continuam in sensu communi, i. e. nullibi interruptam sed ubique cohaerentem.}
